
\InputIfFileExists{ztex_doc-cfg.tex}{}{}
\documentclass[
  hyper=true,
  lang=cn,
  class=l3doc, 
  classOption={11pt},
]{../code/ztex}
\input{ztex_interface_preamble.tex}

\begin{document}
\section{Test-1}
\subsection{Test-2}

\includegraphicsifexsit[height=.95\paperheight]
  {./support/pics/toctree.pdf}
  {example-image-a}

\end{document}



\InputIfFileExists{ztex_doc-cfg.tex}{}{}
% \documentclass[
%   hyper=true,
%   class=l3doc, 
%   classOption={11pt},
% ]{../code/ztex}
% \ztexloadlib{primitive}
% \uncompressPDF

% \begin{document}
% Hello world
% \end{document}
\documentclass[
  lang=cn,                      %
  % hyper=true,
  class=l3doc,                  %
  classOption={11pt},
]{../code/ztex}
% \ztexloadlib{primitive}
% \uncompressPDF
% \input{ztex_interface_preamble.tex}
\geometry{paperwidth=32cm, paperheight=81.5cm, margin=1cm}
% \geometry{paperwidth=25cm, paperheight=28cm, margin=.5cm}
\usepackage{tikz}
\usepackage{tikz-qtree}
\usepackage[log]{tikz-qtree-patch}
% \usetikzlibrary { mindmap, trees } % for example VII
\setcounter{tocdepth}{4}
\ztocenable[toc=ztex_interface, lom, lot, lof]
% \usepackage{etl}



\begin{document}
% % % Example XI: test longtabletoc
% \definecolor{SeaGreen}{rgb}{0.18, 0.55, 0.34}
% \def\customtocthemecolor{SeaGreen}
% \ExplSyntaxOn\makeatletter
% \dim_new:N \linewidthfix
% \def\tocupdatelinewidthfix{
%   \dim_set:Nn \linewidthfix {\linewidth-2\fboxsep}
% }
% \def\ornamentcorner#1#2#3{
%   \node[anchor=#1, rotate=#3, outer~sep=0pt, inner~sep=0pt] at (frame.#2)
%   {\pgfornamenthan[width=.8cm, color=\customtocthemecolor]{10}};}
% \def\ztochyperpage#1{\exp_args:Nne \hyper@link{link}{page.#1}{#1}}
% \makeatother\ExplSyntaxOff

% \begin{dxample}[1]
% \ztocgroupinsert{section,1,before}
%   {\begin{longtable}{|l|l|r|}}
% \ztocgroupinsert{section,4,after}
%   {\\ \hline\end{longtable}}
% {
%   \ztocformat\section{
%     explicit=true,
%     code = {%
%       \ztociffirst{\kill\hline\multicolumn{3}{|c|}{\large\bfseries Table of Contents}\\\hline}{\\\hline}%
%       \multicolumn{3}{||c||}{\bfseries 第\;\ztocname\;节\enskip\ztoctitle}%
%     }
%   }
%   \ztocformat\subsection{
%     explicit=true,
%     code ={\\\hline
%       \textbf{\S\ztocname} & \ztocname\enskip\ztoctitle & \ztochyperpage{\ztocpage}
%     }
%   }
%   \ztocformat\subsubsection{
%     explicit=true,
%     code = {\\
%       & \ztocname\enskip\ztoctitle & \ztochyperpage{\ztocpage}
%     }
%   }
%   \zlocaltoc{section}{1, 2, 7, 9}
% }
% \end{dxample}



% % Example X: test dxample env:
% % \docsourcelinenotostr
% % Example lineno to path:
% \begin{dxample}[0]
% has\hdsah jkasjd
% klasjdf

% \jklajs

% d
% \end{dxample}

% \begin{DocExample}
% has\hdsah jkasjd
% klasjdf

% \jklajs

% d
% \end{DocExample}
% \ztexDocPrintSource



% % Example IX: last toc example from CUS
% \newgeometry{hmargin=3cm, vmargin=1in}
% \ExplSyntaxOn
% \clist_set:Nn \l_tmpa_clist { 0pt, 1.5em, 2.3em, 3.2em }
% \cs_new:Npn \__accum_skip:n #1
%   {
%     \exp_args:Ne \dim_eval:n
%       {
%         \int_step_tokens:nn {#1}
%           { \__get_dim:n }
%         0pt
%       }
%   }
% \cs_new:Npn \__get_dim:n #1
%   { \clist_item:Nn \l_tmpa_clist {#1} }
% % \edef\TTT{\__accum_skip:n {3}}
% % \show\TTT


% \newcommand{\nbraket}
%   { \makebox[0pt]{\rule[-7.5pt]{.75pt}{19pt}} }
% \newcommand{\ubraket}
%   {
%     \makebox[0pt]{
%       \rule[-7pt]{.75pt}{10pt}
%     }
%     \makebox[0pt][l]{
%       \kern-.375pt\raise3pt\hbox{\rule{3pt}{.75pt}}
%     }
%   }
% \newcommand{\bbraket}
%   {
%     \makebox[0pt]{ \rule[3pt]{.75pt}{10pt} }
%     \makebox[0pt][l]
%       { \kern-.375pt\raise3pt
%       \hbox{\rule{3pt}{.75pt}} }
%   }
% \newcommand{\showbar}
%   {
%     \prg_replicate:nn {\ztocdepth - 2}
%       { \kern20pt \nbraket }
%     \kern20pt \ztociffirst { \ubraket }
%       { \ztociflast{\bbraket}{\nbraket} }
%   }
% \ztocformat{\section, \subsection, \subsubsection}
%   {
%     % name.before = \llap{\smash{\showbar}},
%     % space.left=0pt,
%     % leader.content=,
%     explicit=true,
%     code = {
%       \noindent
%       \prg_replicate:nn {\ztocdepth-1}
%         {
%           \hskip 20pt\relax
%         }
%       \llap{\smash{\showbar}}
%       \kern10pt
%       \ztoctitle\par
%     }
%   }
% \ExplSyntaxOff
% \thispagestyle{empty}
% {
%   {
%     \ztocformat\subsection{ignore.name={10.}}
%     \ztocformat\subsubsection{ignore.name={10.}}
%     \multicolcontents{2}
%   }
%   \clearpage
% }
% \restoregeometry


% \begin{function}[added=2025-09-12, pTF, rEXP]{
%     \ztoc_if_last_tocline:Nnn,
%     \ztoc_if_last_tocline:Nee,
%     \ztoc_if_last_tocline:cee,
%   }
%   \begin{syntax}
%   \cmd{\ztoc_if_last_tocline:NnnTF}
%   ~~~~\meta{seq}\Arg{name}\Arg{title}
%   ~~~~\Arg{true code}\Arg{false code}
%   \end{syntax}
%   此命令用于判断当前目录是否为其所在 level 的最后一项, \meta{seq} 用于指定当前所排版目录
%   对应的 keyval seq 变量. \textbf{备注}: 在 \texttt{e}-type 中, 此命令\textbf{不能}被完全展开.
%   \vskip1.25em\relax
%   \noindent{\sffamily\color{red}NOTE: 当最后的目录条目为 
%   ``\verb|\contentsline{subsection}{}{}{}|, \\ \verb|\contentsline{subsection}{}{}{}|'' 时, 
%   此命令对于上述 ``\texttt{section}'' 条目的判断会出错, 此时用户可以手动添加对应的代码.}
% \end{function}


% % >> Example VIII: test table/figure name
% \figurename, \tablename

% \fbox{\zcmdgentoken{123}{12}, \zcmdgentoken{65}{11}}

% % \marginpar{\zsectcaption{figure}{Test Margin Figure}}

% \begin{table}[!htb]
%   \begin{tabular}{ccc}
%     1 & 2 & 3 \\
%     1 & 2 & 3
%   \end{tabular}
%   \caption{Test caption}
% \end{table}

% % \ExplSyntaxOn
% % \edef\TTT
% %   {
% %     \ztex_tl_replace_all:nnn {[ab[c]d}{[}{(}
% %   }
% % \show\TTT
% % % > \TTT=macro:
% % % ->c]d.
% % \ExplSyntaxOff

% \section{\texorpdfstring{\ztex{}}{zTeX} 源码}
% \vskip2em
% {\fboxsep0pt\relax
%   \begin{framed}%
%   \begin{multicols}{2}%
%   \ztocformat\section{ignore=true}%
%   \ztocformat\subsubsection{name.width=0pt, name.after=\space}%
%   \zlocaltoc{section}{10}%
%   \end{multicols}%
%   \end{framed}%
% }

% \textbf{Eureka~\rlap{\wscale*{6em}{\normalfont\small\sffamily\color{gray}zongpingding5(at)outlook(dot)com}}}

% % 也能实现比较复杂的盒子排版\Footnote{用户可以参考 \pkg{longfbox} 宏包的文档}

% % \begin{table}[!htb]
% %   \begin{tblr}{
% %     colspec={|X[1.25, c]|X[1, c]|X[1, c]|X[1, c]|},
% %     rowspec={|Q[m]|Q[m]|Q[m]|Q[m]|Q[m]|},
% %     cells={cmd=\pkg},
% %     note{b} = {XXX-NOTE},
% %   }
% %   geometry  & fancyhdr       & graphicx              & xcolor   \\
% %   amsmath   & amsfonts       & esint                 & etoolbox\TblrNote{b} \\
% %   framed    & framedmulticol & cleveref/zref-clever  &   \\ 
% %   \end{tblr}
% %   \caption{\ztex{} 文档类基本宏包}
% %   \label{tab:basic-package}
% % \end{table}

% % % \footnotetext[2]{加载 \ztex{} 后, 用户可以直接使用 \pkg{ctexpatch} 或 \pkg{xpatch} 宏包提供的命令.}

% % 也能实现比较好可爱的计划\Footnote{工行卡说的话 \pkg{longfbox} 工行卡说的话}


% >> Example VI: test compile time
% \input{ztex_interface.toc}
% \ExplSyntaxOn
% % \ztocformat\section{
% %   explicit=true, enhance=true,
% %   code = { 
% %      (\ztocthehilevel) 
% %      \bool_show:N \l__ztoc_public_macro_enhance_bool
% %      > \l__ztoc_public_macro_enhance_bool=true.
% %    }
% % }
% % \contentsline{section}{{1}{基本介绍}}{1}{section.1}
% % \bool_show:N \l__ztoc_public_macro_enhance_bool
% % > \l__ztoc_public_macro_enhance_bool=false.
% \ExplSyntaxOff
% % \ztochilevel{toc}
% \end{document}



% >> Example VII: mindmap toc
% \ExplSyntaxOn
% \int_new:N \g__mindmap_color_int
% \seq_new:N \g__mindmap_data_seq
% \seq_new:N \g__mindmap_keyvaldata_seq
% \def\toclinecmd#1#2#3#4
%   {
%     {\ztexlink{#1.#2}{\textcolor{white}{#3(#4)}}}
%   }

% \cs_new:Npn \__ztoc_gen_mindmap_data:nn #1#2
%   {
%     \ztoc_table_filter_byclass:nnNN
%       { #1 }{ #2 } 
%       \g_ztoc_keyvaltoc_seq
%       \g__mindmap_data_seq

%     \ztoc_generate_keyvaltable_seq:NN
%       \g__mindmap_data_seq
%       \g__mindmap_keyvaldata_seq
%     % \seq_show:N \g__mindmap_keyvaldata_seq

%     \ztoc_group_parser:NNNnnnnn \g__mindmap_keyvaldata_seq
%       \toclinecmd \g_tmpa_tl 
%       { 
%         child [ concept~color = blue!
%           \int_eval:n{\g__mindmap_color_int*6}!green
%         ] \bgroup node[concept]
%       }
%       {}{}{\egroup}{\int_gincr:N \g__mindmap_color_int}
%     % \tl_show:N \g_tmpa_tl

%     \exp_args:NNo \str_set:Nn \l_tmpa_str { \g_tmpa_tl }
%     \exp_args:NNne \str_replace_all:Nnn \l_tmpa_str
%       { \bgroup }{ \char_generate:nn {123}{12} }
%     \exp_args:NNne \str_replace_all:Nnn \l_tmpa_str
%       { \egroup }{ \char_generate:nn {125}{12} }
%     % \str_show:N \l_tmpa_str
%     \tl_gset_rescan:Nno \g_tmpa_tl
%       { \cctab_select:N \c_document_cctab }
%       { \l_tmpa_str }
%   }
% \cs_new:Npn \__tkz_mind_map:n #1
%   {
%     \node[concept] {\color{white}\ztex{}~Interface}
%       #1;
%   }
% \newcommand{\mindmap}[2]
%   {
%     \__ztoc_gen_mindmap_data:nn
%       { #1 }{ #2 }
%     % \tl_show:N \g_tmpa_tl
%     \exp_args:No \__tkz_mind_map:n
%       { \g_tmpa_tl }
%   }

% % >> this replace fails:
% % \exp_args:NNne \tl_replace_all:Nnn \g_tmpa_tl
% %   { \bgroup }{ \char_generate:nn {123}{1} }
% % \exp_args:NNne \tl_replace_all:Nnn \g_tmpa_tl
% %   { \egroup }{ \char_generate:nn {125}{2} }


% % \tl_show:N \g_tmpa_tl
% % \g_tmpa_tl=child{node[concept]{基本介绍}}child{node[concept]{安装使用}
% % child{node[concept]{在线模板}}child{node[concept]{本地安装}}child{node[concept]{快速开始}}}.

% % \edef\TTT{
% %   % \ztex_tl_replace_all:onn {\l_tmpa_tl} % has bug
% %   % \exp_args:No \etl_replace_all:nnn {\l_tmpa_tl}
% %     { [ }{ ! }
% % }
% % \show\TTT
% \ExplSyntaxOff

% \thispagestyle{empty}
% \begin{tikzpicture}[
%     mindmap, grow cyclic,
%     text width=2cm, align=flush center,
%     nodes={concept}, concept color=blue!60,
%     root concept/.append style={text width=4cm, font=\Large},
%     level 1/.append style={level distance=5cm, sibling angle=45, text width=3cm, font=\Large},
%     level 2/.append style={level distance=9cm, sibling angle=90, text width=3cm, font=\Large},
%     level 3/.append style={level distance=5cm, sibling angle=45, text width=3cm, font=\Large},
%     level 1 concept/.append style={font=\normalsize, color=white},
%   ]
%   % \node[concept] {Computer Science}
%   %   child[concept color=orange]
%   %     {
%   %       node[concept] {Practical}
%   %         child {node[concept] {Pro\-gramming Languages}} 
%   %         child {node[concept] {Software Engineering}}
%   %         child {node[concept] {Data Structures}}
%   %         child {node[concept] {Algorithms}}
%   %     }
%   %   child[concept color=blue!50!black] 
%   %     {
%   %       node[concept] (app) {Applied}
%   %         child {node[concept] {Data Base}}
%   %         child {node[concept] {WWW}}
%   %     };

%   % test example i: ---> does NOT work
%   % \node[concept] {\ztex{} Interface}
%   % child\bgroup node[concept]{基本介绍}\egroup child\bgroup
%   %   node[concept]{安装使用}child\bgroup node[concept]{在线模板}\egroup child\bgroup
%   %   node[concept]{本地安装}\egroup child\bgroup node[concept]{快速开始}\egroup \egroup;

%   % test example ii: ---> works
%   % \node[concept] {\ztex{} Interface}
%   %   child{ node[concept]{基本介绍} } 
%   %   child{ 
%   %     node[concept]{安装使用}
%   %     child{ node[concept]{在线模板} } 
%   %     child{ node[concept]{本地安装} } 
%   %     child{ node[concept]{快速开始} } 
%   %   };
  
%   % test example iii: ---> works
%   % \node[concept] {\ztex{} Interface} child{node[concept]{基本介绍}}child{node[concept]{安装使用}child{node[concept]{在线模板}}child{node[concept]{本地安装}}child{node[concept]{快速开始}}};

%   \mindmap{section}{7}
% \end{tikzpicture}

% \begin{dxample}[1]
% \newgeometry{margin=0pt}
% \mbox{}\vfill
% % \usetikzlibrary { mindmap, trees }
% \ExplSyntaxOn
% \int_new:N \g__mindmap_color_int
% \seq_new:N \g__mindmap_data_seq
% \seq_new:N \g__mindmap_keyvaldata_seq
% \def\toclinecmd#1#2#3#4
%   {
%     % override hyperlink color.
%     {\ztexlink{#1.#2}{\textcolor{white}{#3(#4)}}}
%   }
% \cs_new:Npn \__ztoc_gen_mindmap_data:nn #1#2
%   {
%     \ztoc_table_filter_byclass:nnNN
%       { #1 }{ #2 } 
%       \g_ztoc_keyvaltoc_seq
%       \g__mindmap_data_seq
%     \ztoc_generate_keyvaltable_seq:NN
%       \g__mindmap_data_seq
%       \g__mindmap_keyvaldata_seq
%     \ztoc_group_parser:NNNnnnnn \g__mindmap_keyvaldata_seq
%       \toclinecmd \g_tmpa_tl
%         {
%           child [ concept~color = blue!
%             \int_eval:n{\g__mindmap_color_int*6}!green
%           ] \bgroup node[concept]
%         }
%       {}{}{\egroup}{\int_gincr:N \g__mindmap_color_int}
%     % replace '\bgroup' and '\egroup' by '{' and '}'. 
%     \exp_args:NNo \str_set:Nn \l_tmpa_str { \g_tmpa_tl }
%     \exp_args:NNne \str_replace_all:Nnn \l_tmpa_str
%       { \bgroup }{ \char_generate:nn {123}{12} }
%     \exp_args:NNne \str_replace_all:Nnn \l_tmpa_str
%       { \egroup }{ \char_generate:nn {125}{12} }
%     \tl_gset_rescan:Nno \g_tmpa_tl
%       { \cctab_select:N \c_document_cctab }
%       { \l_tmpa_str }
%   }
% \cs_new:Npn \__tkz_mind_map:n #1
%   {
%     \tl_if_empty:nF {#1}
%       {
%         \node[concept] 
%           {\color{white}\ztex{}~Interface}
%           #1;
%       }
%   }
% \newcommand{\mindmap}[2]
%   {
%     \__ztoc_gen_mindmap_data:nn
%       { #1 }{ #2 }
%     \exp_args:No \__tkz_mind_map:n
%       { \g_tmpa_tl }
%   }
% \ExplSyntaxOff
% \noindent\begin{tikzpicture}[
%     mindmap, grow cyclic,
%     text width=2cm, align=flush center,
%     nodes={concept}, concept color=blue!60,
%     root concept/.append style={text width=4cm, font=\Large},
%     level 1/.append style={level distance=5cm, sibling angle=45, text width=3cm, font=\Large},
%     level 2/.append style={level distance=8.75cm, sibling angle=90, text width=3cm, font=\Large},
%     level 3/.append style={level distance=5cm, sibling angle=45, text width=3cm, font=\Large},
%     level 1 concept/.append style={font=\normalsize, color=white},
%   ]
%   \mindmap{section}{7}
% \end{tikzpicture}
% \vfill\mbox{}
% % \end{dxample}
% \restoregeometry



% >> example I:
\definecolor{Red}{HTML}{e32b2d}
\definecolor{Green}{HTML}{15821b}
\definecolor{Blue}{HTML}{0984e3}
\tikzset{
  section/.style={draw, rectangle, rounded corners, fill=Red},
  subsection/.style={draw, rectangle, rounded corners, fill=Blue, text=white},
  subsubsection/.style={draw, rectangle, rounded corners, fill=Green, text=white},
}
\ExplSyntaxOn
% \def\toclinecmd#1#2#3#4{{#1:#2:#3:#4}~}
\def\toclinecmd#1#2#3#4{{\gparserclass(\gparserindex):#2:#3:#4}~}
\ztoc_group_parser:NNNnnnnn \g_ztoc_keyvaltoc_seq
  \toclinecmd \g_tmpa_tl {[.}{}{}{~]}{}

% \tl_show:N \g_tmpa_tl
\def\TocTree
  {
    \exp_args:Nno \__toc_tree:nn
      {\ztex{}~接口文档}
      {\g_tmpa_tl}
  }
\cs_new:Npn \__toc_tree:nn #1#2
  {
    \quark_if_no_value:nF {#1}
      {
        \Tree[.\node[section]{#1};~#2 ]
      }
  }
\ExplSyntaxOff

\thispagestyle{empty}
\begin{tikzpicture}[
    grow'=right, level distance=1cm,
    sibling distance=10pt,
    every tree node/.style = {
      anchor=west, outer sep=0pt, draw, rectangle, 
      rounded corners, fill=gray!25
    },
    every level 1 node/.style = {
      subsection
    },
    every level 2 node/.style = {
      subsubsection
    },
    edge from parent/.style={
      draw, thick, rounded corners,
      edge from parent path={
        (\tikzparentnode.east) -- +(.5cm, 0pt)
        |- (\tikzchildnode.west)
      },
    },
  ]
  \TocTree
\end{tikzpicture}


% >> example II:
% \makeatletter
% \def\ztochyperpage#1{\hyper@link{link}{page.#1}{#1}}
% \makeatother
% \ztocgroupinsert{section,1,before}
%   {\begin{longtable}{|l|l|r|}}
% \ztocgroupinsert{section,4,after}
%   {\\ \hline\end{longtable}}
% {
%   \ztocformat\section{
%     explicit=true,
%     code = {%
%       \ztociffirst{\kill\hline\multicolumn{3}{|c|}{\large\bfseries Table of Contents}\\\hline}{\\\hline}%
%       \multicolumn{3}{||c||}{\bfseries 第\;\ztocname\;节\enskip\ztoctitle}%
%     }
%   }
%   \ztocformat\subsection{
%     explicit=true,
%     code ={\\\hline
%       \textbf{\S\ztocname} & \ztocname\enskip\ztoctitle & \ztochyperpage{\ztocpage}
%     }
%   }
%   \ztocformat\subsubsection{
%     explicit=true,
%     code = {\\
%       & \ztocname\enskip\ztoctitle & \ztochyperpage{\ztocpage}
%     }
%   }
%   \zlocaltoc{section}{1, 2, 7, 9}
% }


% >> example: III
% \setcounter{tocdepth}{3}
% \begingroup\makeatletter
% \definecolor{SeaGreen}{rgb}{0.18, 0.55, 0.34}
% \def\ornamentcorner#1#2#3{%
%   \node[anchor=#1, rotate=#3, outer sep=0pt, inner sep=0pt] at (frame.#2)
%     {\pgfornamenthan[width=.8cm, color=\customtocthemecolor]{10}};}
% % \setcounter{tocdepth}{3}
% \def\customtocthemecolor{SeaGreen}
% \ztocformat\section
%   {
%     explicit=true,
%     code={%
%       \noindent\hskip1.5em\fcolorbox{white}{\customtocthemecolor}{\hb@xt@
%         \dimexpr\linewidth-2\fboxsep-1.5em\relax
%         {\bfseries\large\color{white}第\zhnumber{\ztocname}节
%           \hskip.75em\relax\ztoctitle\hfill\ztocpage}}%
%     }
%   }
% \ztocformat\subsection{
%   name.before=\S\;, name.width=2.7em,
%   format+=\color{\customtocthemecolor},
%   leader.sep=2pt, page.width=1em,
%   hyper.title=true,
%   leader.content=\textcolor{\customtocthemecolor}{.},
% }
% \ztocgroupinsert{section,1,before}%
%   {%
%     \begin{tcolorbox}[
%       enhanced, sharp corners,
%       boxrule=.5pt, colback=white,
%       left=20pt, colframe=\customtocthemecolor,
%       overlay={
%         \ornamentcorner{north west}{north west}{0}
%         \ornamentcorner{north west}{south west}{90}
%         \ornamentcorner{north west}{south east}{180}
%         \ornamentcorner{north west}{north east}{270}
%       }
%     ]
%     \pgfornament[width = 2.5cm,color = \customtocthemecolor]{10}%
%     \begin{minipage}{.77\linewidth}%
%   }
% \ztocgroupinsert{section,1,begin}{%
%   % \setlength{\columnsep}{0em}
%   \begin{multicols}{2}}
% \ztocgroupinsert{section,1,after}%
%   {%
%     \end{multicols}%
%     \end{minipage}%
%     \end{tcolorbox}%
%   }
% \zlocaltoc{section}{7}
% \makeatother\endgroup


% example V:
% \ztocformat\subsection{ignore.name={10.}}
% \ztocformat\subsubsection{ignore.name={10.}}
% \multicolcontents{2}
\end{document}




\InputIfFileExists{ztex_doc-cfg.tex}{}{}
\documentclass[
  lang=cn, 
  hyper=true,
  class=l3doc, 
  classOption={11pt},
]{../code/ztex}
\input{ztex_interface_preamble.tex}
\usepackage{tikz-qtree}


\ztocenable[toc=ztex_interface, lom, lot, lof]
\title{\zTeX{} 接口文档}
\author{Eureka}
\date{\textsf{\today}}
\begin{document}

% \begin{dxample}*
% % \ztocformat{\subsection, \subsubsection}{hyper.title=true}
% \definecolor{Green}{HTML}{15821b}
% \tikzset{
%   rootKey/.style={draw, rectangle, rounded corners, fill=Green, text=white},
% }

% \ExplSyntaxOn
% \cs_new:Npn \__toc_tree:nn #1#2
%   {
%     \Tree[.\node[rootKey]{#1};~#2 ]
%   }
% \cs_generate_variant:Nn \__toc_tree:nn {oe}
% \def\wraper#1{[.{\exp_not:n {#1}}~]}
% \cs_new:Npn \ztoc_typeset_toctree:nnn #1#2#3
%   {
%     \ztoc_table_filter_key_byclass:nnnNN {#1}{#2}{#3}
%       \g_ztoc_keyvaltoc_seq \g_tmpb_seq
%     \seq_gpop_left:NN \g_tmpb_seq \secRoot
%     \__toc_tree:oe { \secRoot }
%       { \seq_map_function:NN \g_tmpb_seq \wraper }
%   }
% \def\TocTree#1#2#3
%   {
%     \ztoc_typeset_toctree:nnn
%       {#1}{#2}{#3}
%   }
% \ExplSyntaxOff

% \begin{tikzpicture}[
%     grow'=right,
%     level distance=50pt,
%     sibling distance=10pt,
%     level 1/.style={level distance=25pt},
%     every tree node/.style={anchor=west, draw, rectangle, rounded corners, fill=gray!25},
%     every level 0 node/.style={anchor=east},
%     edge from parent/.style={
%       draw, thick, rounded corners,
%       edge from parent path={
%         (\tikzparentnode.east) -- +(10pt, 0pt)
%         |- (\tikzchildnode.west)
%       }
%     },
%   ]
%   \TocTree{subsection}{4}{page}
%   \begin{scope}[xshift=4.5cm]
%     \TocTree{subsection}{11}{title}
%   \end{scope}
% \end{tikzpicture}
% \end{dxample}

% \begin{dxample}*
% \textbf{SSS}
% \def\TTTa{TTTa}

% \end{dxample}
% \show\TTTa % --> solved by '\aftergroup'


\begin{dxample}*
\definecolor{SeaGreen}{rgb}{0.18, 0.55, 0.34}
\def\customtocthemecolor{SeaGreen}
\ExplSyntaxOn\makeatletter
\dim_new:N \linewidthfix
\def\tocupdatelinewidthfix{
  \dim_set:Nn \linewidthfix {\linewidth-2\fboxsep}
}
\def\ornamentcorner#1#2#3{
  \node[anchor=#1, rotate=#3, outer~sep=0pt, inner~sep=0pt] at (frame.#2)
  {\pgfornamenthan[width=.8cm, color=\customtocthemecolor]{10}};}
% \def\ztochyperpage#1{\exp_args:Nne \hyper@link{link}{page.#1}{#1}}
\makeatother\ExplSyntaxOff
\end{dxample}


\begin{DocExample}*
\makeatletter
\def\customssstyle#1{%
  \ztocgroupinsert{subsection,#1,begin}
    {\hangindent2.3em\relax\noindent\hskip2.3em\relax【}
  \ztocgroupinsert{subsection,#1,end}{】\par}
}
\makeatother
\customssstyle{1} \customssstyle{2}
\customssstyle{3} \customssstyle{4}
{
  \ztocformat\subsubsection{
    explicit=true,
    code={%
      \textit{\ztocname\enskip\ztoctitle\enskip
      (\ztochyperpage{\ztocpage});\enskip}}
  }
  \ztocformat\subsection{
    explicit=true,
    code = {%
      \S\;\ztocname\enskip\ztoctitle
      \hskip1em\ztoclink{\ztocpage}\par}
  }
  \begin{multicols}{2}
    \zlocaltoc{subsection}{4, 11, 8, 6}
  \end{multicols}
}
\end{DocExample}
\end{document}


% \ztocgroupinsert{section,1,before}
%   {\begin{longtable}{|l|l|r|}}
% \ztocgroupinsert{section,4,after}
%   {\\ \hline\end{longtable}}
% {
%   \ztocformat\section{
%     explicit=true,
%     code = {%
%       \ztociffirst{\kill\hline\multicolumn{3}{|c|}{\bfseries Table of Contents}\\\hline}{\\\hline}%
%       \multicolumn{3}{||c||}{\bfseries 第\;\ztocname\;节\enskip\ztoctitle}%
%     }
%   }
%   \ztocformat\subsection{
%     explicit=true,
%     code ={\\\hline
%       \textbf{\S\ztocname} & \ztocname\enskip\ztoctitle & \ztocpage
%     }
%   }
%   \ztocformat\subsubsection{
%     explicit=true,
%     code = {\\
%       & \ztocname\enskip\ztoctitle & \ztocpage
%     }
%   }
%   \zlocaltoc[ztex_interface]{section}{1, 2, 7, 9}
% }


% {
%   \ExplSyntaxOn
%   \seq_gclear_new:N \g__ztoc_sec_seq
%   \seq_gclear_new:N \g__ztoc_subsec_seq
%   \newcommand{\saveSectoc}[2]
%     {
%       \seq_if_exist:cF { g__ztoc_sec_#1_seq }
%         { \seq_new:c { g__ztoc_sec_#1_seq } }
%       \exp_args:Nno \seq_gput_right:cn { g__ztoc_sec_#1_seq }
%         { #2 }
%     }
%   \newcommand{\saveSubSectoc}[2]
%     {
%       \seq_if_exist:cF { g__ztoc_sec_#1_child_seq }
%         { \seq_new:c { g__ztoc_sec_#1_child_seq } }
%       \exp_args:Nne \seq_gput_right:cn { g__ztoc_sec_#1_child_seq }
%         { #2 }
%       \seq_show:c { g__ztoc_sec_#1_child_seq }
%     }
%   % % \ztoc_table_filter_bynametitle:nnNNN {3}{基本命令}
%   % \ztoc_table_filter_bynametitle:nnNNN {7.1}{\pkg        {font}~模块} 
%   % % ---> still found the contents line.
%   % % ---> NOTE: all spaces after cs will be ignored when comparing,
%   % %            but other spaces can NOT be removed.
%   %   \g_ztoc_keyvaltoc_seq \l_tmpa_prop \l_tmpa_int
%   % \prop_show:N \l_tmpa_prop
%   % \int_show:N \l_tmpa_int
%   % % case 1: '3' + '基本命令':
%   % % >  {class}  =>  {section}
%   % % >  {name}  =>  {3}
%   % % >  {title}  =>  {基本命令}
%   % % >  {page}  =>  {5}
%   % % >  {raw}  =>  {\contentsline {section}{{3}{基本命令}}{5}{section.3}}.
%   % % > \l_tmpa_int=6. 
%   % % case 2: '7.1' + '\pkg  {font} 模块':
%   % % >  {class}  =>  {subsection}
%   % % >  {name}  =>  {7.1}
%   % % >  {title}  =>  {\pkg {font} 模块}
%   % % >  {page}  =>  {15}
%   % % >  {raw}  =>  {\contentsline {subsection}{{7.1}{\pkg {font}
%   % % 模块}}{15}{subsection.7.1}}.
%   % % > \l_tmpa_int=11.

%   % % test if last contents line:
%   % \ztoc_if_last_tocline:NnnTF \g_ztoc_keyvaltoc_seq
%   %   % {3}{基本命令} % ------>NO
%   %   % {12}{索引} % ------>YES
%   %   % {11.4.5}{primitive} % ------>YES
%   %   {8.4}{\pkg  {thm}~库} % ------>YES
%   %   { \typeout{------>YES} }
%   %   { \typeout{------>NO} }
%   % % NOTE: move the following code out of expl3 to debug:
%   % % \ztocformat\subsection{
%   % %   explicit=true,
%   % %   code ={%
%   % %     \ztocname:\ztoctitle:\ztocpage
%   % %     \ztociflast{LAST}{NO}\par % ---> '\ztociflast' works as expected!
%   % %   }
%   % % }
%   % % Test last line check example II.
%   % % \ztocformat{\subsection, \subsubsection}{
%   % %   explicit=true,
%   % %   code ={%
%   % %     \ztocname:\ztoctitle:\ztocpage
%   % %     \ztociflast{LAST}{NO}\par % ---> '\ztociflast' works as expected!
%   % %   }
%   % % }
%   % % some bug lies in last line check:
%   % % modified toc file:
%   % %   \contentsline{subsection}{{7.9}{\pkg  {sclist} 模块}}{122}{subsection.7.9}%
%   % %   \contentsline{subsubsection}{{7.9.1}{\pkg  {sclist} 模块-KKK}}{122}{subsection.7.9}%
%   % % PDF Output:
%   % %   7.9:sclist 模块:122NO
%   %    7.9.1:sclist 模块-KKK:122LAST
%   % % NOTE: currently, users need insert the code for last 
%   %         contentsline manually in this (bug)edeg case .
%   \ExplSyntaxOff
%   % \ztocformat{\subsubsection, \subsection}{title.after=(\ztociflast{LAST}{NO})}
%   %% wraper toc in tikz mindmap:
%   \ztocformat\section
%     {
%       explicit=true,
%       code={%
%         \saveSectoc{\ztocname}{\ztoctitle\enskip(\ztocpage)}%
%       }
%     }
%   \ztocformat\subsection
%     {
%       explicit=true,
%       code={%
%         \saveSubSectoc{\ztocname}{\ztoctitle\enskip(\ztocpage)}%
%       }
%     }
%   \zlocaltoc[ztex_interface]{section}{1, 2, 3}
% }

% {
% \ExplSyntaxOn
% % \seq_show:N \g_ztoc_keyvaltoc_seq
% % \ztoc_table_filter_byclass:nnNN {subsection}{4}
% %   \g_ztoc_keyvaltoc_seq\g_tmpa_seq
% % % \seq_show:N \g_tmpa_seq

% % \ztoc_table_filter_key_byclass:nnnNN {section}{2}{name}
% %   \g_ztoc_keyvaltoc_seq\g_tmpa_seq
% \ztoc_table_filter_key_byclass:nnnNN {subsection}{4}{title}
%   \g_ztoc_keyvaltoc_seq\g_tmpb_seq

% % \seq_show:N \g_tmpa_seq
% % \seq_show:N \g_tmpb_seq

% \seq_gpop_left:NN \g_tmpb_seq \secRoot
% % \show\secRoot
% % > \secRoot=macro:
% % ->\pkg {font} 模块.
% \def\wraper#1{[.{\exp_not:n {#1}}~]}
% \edef\TTTa{\seq_map_function:NN \g_tmpb_seq \wraper }
% \cs_new:Npn \__toc_tree:nn #1#2
%   {
%     % \def\TTTx{#1}
%     % \def\TTTy{#2}
%     % \show\TTTx
%     % \show\TTTy
%     \Tree[.{#1}~
%       #2
%     ]
%   }
% \cs_generate_variant:Nn \__toc_tree:nn {oe}
% % \exp_args:NNo \def\TocTree{\exp_not:N \Tree[.{\secRoot}~\TTTa]}
% \gdef\TocTree#1#2{\__toc_tree:oe {#1}{#2}}
% % \TocTree{\secRoot}{\TTTa}
% % \show\TocTree
% % \show\TTTa
% % > \TTTa=macro:
% % ->[.{\pkg {font} 模块} ][.{字体机制} ][.{默认字体族} ][.{新建字体族} ]
% % [.{切换字体} ][.{\ztex {} 接口}][.{杂项} ].



% % \ztoc_table_filter_key_byclass:nnnNN {subsection}{4}{title}
% %   \g_ztoc_keyvaltoc_seq\g_tmpb_seq
% % \seq_gpop_left:NN \g_tmpb_seq \secRoot
% % \edef\TTTa{\seq_map_function:NN \g_tmpb_seq \wraper }
% % \cs_generate_variant:Nn \__toc_tree:nn {oe}
% % \gdef\TocTree#1#2{\__toc_tree:oe {#1}{#2}}
% \ExplSyntaxOff

% \begin{tikzpicture}[
%     grow'=right,
%     level distance=50pt,
%     sibling distance=10pt,
%     level 1/.style={level distance=25pt},
%     every tree node/.style={anchor=west, draw, rectangle, rounded corners, fill=gray!25},
%     every level 0 node/.style={anchor=east},
%     edge from parent/.style={
%       draw, thick, rounded corners,
%       edge from parent path={
%         (\tikzparentnode.east) -- +(10pt, 0pt)
%         |- (\tikzchildnode.west)
%       }
%     },
%   ]
%   \TocTree{\secRoot}{\TTTa}
%   % \Tree
%   %   [%.\node[boolKey] {Euclidean Space $\mathbb{R}^n$};
%   %     % [.{Inner Product: $\langle x, y\rangle = \langle y, x\rangle$} ]
%   %     % [.{Fatou's Lemma: $\displaystyle\int\liminf_{k\to\infty} f_k dx \le \liminf_{k\to\infty} \int f_k dx$} ]
%   %     % [.{Functional Analysis} ]
%   %     % [.\node[metaKey] {\parbox{11em}{A \textbf{Vitali} set is an elementary example of a 
%   %     %   set of real numbers that is not \textbf{Lebesgue measurable}, found by Giuseppe 
%   %     %   Vitali in 1905.}}; ]
%   %     % [.Vector
%   %     %   [.{Triangle Inequality: $\|x+y\| \le \|x\| + \|y\|$} ]
%   %     %   [.BBB-5 ]
%   %     % ]
%   %     .\node[boolKey] {\pkg {font} 模块};
%   %     [.{字体机制} ]
%   %     [.{默认字体族} ][.{新建字体族} ]
%   %     [.{切换字体} ][.{\ztex {} 接口} ]
%   %     [.{杂项} ]
%   %   ]
% \end{tikzpicture}

% }
